% =============================================================================
% ISI Country Brief — Estonia
% External Supplier Concentration Profile
%
% ISI Paper Series — Country Brief
% International Sovereignty Institute — Research Unit
% March 2026
% =============================================================================
\documentclass[12pt,a4paper]{article}

% --- ISI Paper Series shared template ----------------------------------------
\usepackage{isi-series}

% --- PDF metadata ------------------------------------------------------------
\hypersetup{
  pdftitle={External Supplier Concentration Profile: Estonia — ISI EU-27 2024},
  pdfauthor={International Sovereignty Institute -- Research Unit},
  pdfcreator={International Sovereignty Institute},
  pdfsubject={ISI Paper Series -- Country Brief EE, EU-27, Vintage 2024},
  pdfkeywords={sovereignty index, EU-27, supplier concentration,
    estonia, country brief},
}

% ==============================================================================
\begin{document}
\hypersetup{pageanchor=false}

% ---------- Title Page with Metadata ------------------------------------------
\begin{titlepage}
\centering
\vspace*{2cm}
{\Large\textsc{ISI Paper Series --- Country Brief}}\\[1.2cm]
{\LARGE\bfseries External Supplier Concentration Profile:\\[0.3em]
Estonia}\\[1.5cm]
{\large International Sovereignty Index (ISI)\\[0.3em]
EU-27 --- Vintage 2024 --- Methodology v1.0}\\[1.2cm]
{\large International Sovereignty Institute --- Research Unit}\\[0.8cm]
{\large March 2026}

\vspace{0.8cm}
\noindent\rule{\textwidth}{0.4pt}
\vspace{0.4cm}

% --- Research Metadata --------------------------------------------------------
\begin{flushleft}
\noindent\textbf{Data basis:} ISI Paper~2 --- EU-27 Empirical Results 2024 (Methodology~v1.0, Vintage~2024).\\[4pt]
\noindent\textbf{Methodology:} ISI Paper~1 --- Methodological Foundations, Version~1.0.

\medskip

\noindent\textbf{JEL Codes:} F02, F15, F52, O30, Q43\\[4pt]
\noindent\textbf{Keywords:} sovereignty index; EU-27; supplier concentration; composite indicator; variance decomposition; defence dependence

\medskip

\noindent\textbf{Related publications in the ISI Paper Series:}
\begin{itemize}[nosep,leftmargin=1.5em]
\item ISI Paper~1 --- Technical Specification (v1.0) --- doi:\\
  \url{https://doi.org/10.5281/zenodo.18764227}
\item ISI Paper~2 --- EU-27 Empirical Results 2024 (v1.0) --- doi:\\
  \url{https://doi.org/10.5281/zenodo.18764170}
\end{itemize}
\end{flushleft}
\vfill
\end{titlepage}

\hypersetup{pageanchor=true}
\pagenumbering{arabic}
\pagestyle{fancy}
\fancyhead[L]{\small\sffamily\textit{ISI Country Brief}}
\fancyhead[R]{\small\sffamily\textit{Estonia}}
\renewcommand{\headrulewidth}{0.4pt}

% ---------- Body --------------------------------------------------------------

\section{Executive Summary}
\label{sec:summary}

Estonia ranks~13th in the EU-27 with a composite \ISI score of~0.343, classified as \emph{moderately concentrated}.  The dominant axes are logistics~(0.470) and energy~(0.439).  The structural profile is characterised as dual-axis concentration (logistics and energy).  Under Dirichlet weight perturbation (10\,000~Monte Carlo draws), Estonia's mean rank is~13.46 with a standard deviation of~2.21; the 5th--95th percentile band spans ranks~9 to~17.


\section{Country Position in the EU-27}
\label{sec:position}

Estonia occupies rank~13 of~27 EU member states (composite score~0.343).  The EU-27 mean composite score is~0.344 with a standard deviation of~0.070.  Estonia's score lies below the EU-27 mean, indicating below-average external supplier concentration.  In the ranking, Estonia is situated between Luxembourg (rank~12; 0.360) and Bulgaria (rank~14; 0.340).


\section{Axis Profile}
\label{sec:axis-profile}

\Cref{tab:axis-profile} presents Estonia's scores on all six \ISI axes.

\begin{table}[htbp]
\centering\small
\caption{Estonia's \ISI axis profile.}
\label{tab:axis-profile}
\begin{tabular}{@{}lr@{}}
\toprule
\textbf{Axis} & \textbf{Score} \\
\midrule
    Finance           & 0.182 \\
    Energy            & 0.439 \\
    Technology        & 0.227 \\
    Defence           & 0.356 \\
    Critical Inputs   & 0.386 \\
    Logistics         & 0.470 \\
\bottomrule
\end{tabular}
\end{table}

The highest axis score is logistics~(0.470), which lies below the EU-27 axis mean of~0.518.  The second-highest axis is energy~(0.439), above the EU-27 axis mean of~0.365.  The lowest axis score is finance~(0.182), above the EU-27 mean of~0.148.


\section{Structural Drivers}
\label{sec:drivers}

\begin{sloppypar}
Estonia's composite score is driven by two axes in combination: logistics~(0.470) and energy~(0.439).  The gap between these two axes is~0.031, indicating a dual-axis concentration profile.  Neither axis alone dominates the composite; rather, the combination of elevated scores on both dimensions determines Estonia's position in the ranking.
\end{sloppypar}


\section{Comparative Context}
\label{sec:comparative}

Across the EU-27, the covariance-based variance decomposition shows that defence (35.1\,\%) and logistics (34.3\,\%) together account for 69.4\,\% of total cross-sectional composite variance.  Estonia's defence score~(0.356) lies below the EU-27 defence mean~(0.573), and its logistics score~(0.470) lies below the EU-27 logistics mean~(0.518).  Under Dirichlet weight perturbation, Estonia's rank standard deviation is~2.21, indicating elevated rank volatility relative to the EU-27 median.



\begin{sloppypar}
Estonia's axis profile is comparatively balanced: the spread between the highest axis---logistics~(0.470)---and the lowest---finance~(0.182)---is~0.288, one of the smaller intra-profile ranges in the EU-27.  No single axis towers above the others, and four of the six axes lie within a narrow band.  This balanced structure contrasts with the asymmetric profiles of countries ranked immediately above and below.
\end{sloppypar}

The consequence for rank stability is ambiguous.  On the one hand, balanced profiles are less sensitive to the reweighting of any individual axis.  On the other, Estonia's scores on the two high-variance axes---defence~(0.356) and logistics~(0.470)---are both below the EU-27 means, meaning that the baseline weighting does not amplify the country's strongest dimensions.  The rank standard deviation of~2.21 and the percentile band~(9--17) reflect this: Estonia's position shifts meaningfully depending on whether the weighting favours the high-variance axes or distributes weight more evenly.

Estonia's composite of~0.343 is close to the EU-27 mean~(0.344), placing it near the centre of the cross-sectional distribution---a position where small absolute changes translate into large ordinal movements.


\section{Interpretation Boundary}
\label{sec:interpretation}

The \ISI measures concentration of external suppliers, not sovereignty in a normative or political sense.  High concentration may reflect geography, economic specialisation, or alliance structures.  A high composite score does not imply policy failure; a low score does not imply policy success.  The index provides an empirical measurement basis --- what policymakers derive from it is a separate question.


% ---------- References --------------------------------------------------------
\clearpage
\vspace*{1.5cm}

\begingroup\emergencystretch=2em\hbadness=10000
\section*{References}

\begin{enumerate}[label={[\arabic*]},leftmargin=2em,itemsep=6pt]
\item International Sovereignty Institute.
  \textit{International Sovereignty Index (ISI) technical specification, version~1.0}.
  Zenodo, 2026.\\ISI Paper Series, Paper~1.
  doi:\,\href{https://doi.org/10.5281/zenodo.18764227}{10.5281/zenodo.18764227}.

\item International Sovereignty Institute.
  \textit{International Sovereignty Index (ISI): EU-27 empirical results 2024 (v1.0)}.
  Zenodo, 2026.\\ISI Paper Series, Paper~2.
  doi:\,\href{https://doi.org/10.5281/zenodo.18764170}{10.5281/zenodo.18764170}.
\end{enumerate}
\endgroup

\end{document}
